% =============================================================================
\section{Introduction}
\label{sec:intro}
% =============================================================================

Large Language Models (LLMs) have demonstrated remarkable capabilities as zero-shot passage rerankers in information retrieval (IR) systems~\citep{sun2023chatgpt, qin2024large, zhuang2024setwise}. Given a query and a set of candidate documents retrieved by a first-stage retriever, LLMs can reorder candidates by relevance without any task-specific fine-tuning. Four prompting paradigms have been established: \emph{pointwise} (score each document independently), \emph{pairwise} (compare documents in pairs), \emph{listwise} (rerank a full list), and \emph{setwise} (select the best from a small candidate set)~\citep{zhuang2024setwise}.

The \textbf{setwise approach}~\citep{zhuang2024setwise} occupies a sweet spot in the effectiveness--efficiency trade-off. By presenting a small set of candidates (e.g., 4 documents) and asking the LLM ``\emph{which passage is the most relevant?}'', it extracts more comparative information than pairwise (which compares only 2) while requiring shorter prompts than listwise (which presents 10--20). Combined with classical sorting algorithms like heapsort and bubblesort, setwise ranking achieves high effectiveness with significantly fewer LLM calls than pairwise approaches.

However, current setwise methods have a fundamental limitation: \textbf{each LLM call extracts only one ranking decision}---the identity of the \emph{most relevant} document. The comparative assessment the LLM performs internally necessarily evaluates all candidates, yet we discard all information except the winner. This is analogous to running a tournament but only recording gold medals while throwing away all other placement information.

Consider a concrete example: when an LLM is presented with four passages and asked ``which is the most relevant?'', it must internally assess all four before producing an answer. The model effectively \emph{knows} which passage is the worst too---it simply is not asked to report it. From an information-theoretic perspective, selecting the best from $m$ candidates yields $\log_2(m)$ bits of information. But by asking for \emph{both} the best and worst, we obtain $\log_2(m) + \log_2(m-1)$ bits---nearly doubling the information extracted per call with negligible prompt overhead.

In this paper, we ask: \emph{can we extract more ranking information from each LLM comparison?} We investigate three bidirectional strategies:

\begin{enumerate}[leftmargin=*]
\item \textbf{Bottom-Up Setwise Ranking (\bottomup{})}: We reverse the prompt to ask ``\emph{which passage is the least relevant?}'' and build the ranking from the bottom up by iteratively eliminating the worst documents. This tests whether LLMs exhibit different accuracy and bias patterns when identifying irrelevant versus relevant documents.

\item \textbf{Dual-End Setwise Ranking (\dualend{})}: We modify the prompt to ask for \emph{both} the most and least relevant documents simultaneously: ``\emph{which is the most relevant and which is the least relevant?}'' This extracts two ranking decisions per LLM call. We pair this prompt with two new sorting algorithms---cocktail shaker sort and double-ended selection sort---that are designed to consume both outputs.

\item \textbf{Bidirectional Ensemble (\bidir{})}: We run both top-down (standard) and bottom-up ranking independently and fuse the results using rank aggregation methods. This tests whether the two directions provide complementary relevance signals.
\end{enumerate}

Figure~\ref{fig:overview} illustrates the key differences between standard top-down, bottom-up, and dual-end setwise ranking.

\begin{figure}[!htbp]
\centering
\begin{tikzpicture}[
    doc/.style={rectangle, draw, minimum width=0.6cm, minimum height=0.5cm, font=\scriptsize},
    arrow/.style={-{Stealth[length=2mm]}, thick},
    label/.style={font=\small\bfseries},
    >=Stealth
]
% TopDown
\node[label] at (0, 2.6) {\topdown{}};
\node[doc, fill=green!30] (ta) at (-0.9, 2.0) {A};
\node[doc] (tb) at (-0.3, 2.0) {B};
\node[doc] (tc) at (0.3, 2.0) {C};
\node[doc] (td) at (0.9, 2.0) {D};
\node[font=\scriptsize, text=green!50!black] at (0, 1.3) {``Which is \textbf{best}?'' $\to$ A};
\draw[arrow, green!50!black] (0, 1.1) -- (0, 0.7);
\node[doc, fill=green!30] at (-0.45, 0.4) {\scriptsize 1st};
\node[doc, fill=gray!15] at (0.15, 0.4) {\scriptsize ?};
\node[doc, fill=gray!15] at (0.75, 0.4) {\scriptsize ?};

% BottomUp
\node[label] at (3.4, 2.6) {\bottomup{}};
\node[doc] (ba) at (2.5, 2.0) {A};
\node[doc] (bb) at (3.1, 2.0) {B};
\node[doc] (bc) at (3.7, 2.0) {C};
\node[doc, fill=red!30] (bd) at (4.3, 2.0) {D};
\node[font=\scriptsize, text=red!50!black] at (3.4, 1.3) {``Which is \textbf{worst}?'' $\to$ D};
\draw[arrow, red!50!black] (3.4, 1.1) -- (3.4, 0.7);
\node[doc, fill=gray!15] at (2.65, 0.4) {\scriptsize ?};
\node[doc, fill=gray!15] at (3.25, 0.4) {\scriptsize ?};
\node[doc, fill=red!30] at (3.85, 0.4) {\scriptsize last};

% DualEnd
\node[label] at (6.8, 2.6) {\dualend{}};
\node[doc, fill=green!30] (da) at (5.9, 2.0) {A};
\node[doc] (db) at (6.5, 2.0) {B};
\node[doc] (dc) at (7.1, 2.0) {C};
\node[doc, fill=red!30] (dd) at (7.7, 2.0) {D};
\node[font=\scriptsize, text=blue!50!black, align=center] at (6.8, 1.25) {``\textbf{Best} \& \textbf{worst}?''\\$\to$ A, D};
\draw[arrow, blue!50!black] (6.8, 0.75) -- (6.8, 0.55);
\node[doc, fill=green!30] at (6.05, 0.25) {\scriptsize 1st};
\node[doc, fill=gray!15] at (6.65, 0.25) {\scriptsize ?};
\node[doc, fill=red!30] at (7.25, 0.25) {\scriptsize last};
\end{tikzpicture}
\caption{Comparison of setwise ranking strategies. \topdown{} identifies only the best document per call. \bottomup{} identifies only the worst. \dualend{} identifies both, extracting $\sim$$1.79\times$ ranking information per LLM call.}
\label{fig:overview}
\end{figure}

\textbf{Findings.} Across nine LLMs (Flan-T5 780M--11B, Qwen3 4B--14B, Qwen3.5 4B--27B) and the TREC Deep Learning 2019/2020 benchmarks, we observe a clearer story than the original ``three equally promising directions'' framing would suggest. The headline result is one of \emph{directional asymmetry}: worst-selection alone is not a mirror image of best-selection. Specifically:

\begin{itemize}[leftmargin=*]
\item The \dualend{} family is the strongest overall. \dualend{}-Cocktail achieves the best mean NDCG@10 across all 9 models on both DL19 and DL20, and DualEnd methods (Cocktail or Selection) achieve the highest single-method NDCG@10 in 14 of 18 model--dataset configurations (\S\ref{sec:rq2}).
\item Standalone \bottomup{} is consistently and sometimes catastrophically weaker than \topdown{}. Out of 18 configurations, paired approximate-randomization tests with Bonferroni correction yield \textbf{zero} \bottomup{} wins and \textbf{6} significant losses, including a $-0.23$ NDCG@10 drop on Flan-T5-Large DL19 (\S\ref{sec:rq1}).
\item \bidir{} ensemble does \emph{not} outperform the \topdown{} baseline. The ensemble inherits \bottomup{}'s noise rather than gaining complementary signal: 0 Bonferroni-significant wins and 3 Bonferroni-significant losses, with the optimal fusion weight $\alpha$ trending toward 1.0 (pure TopDown) (\S\ref{sec:rq3}).
\item DualEnd's gains over the strongest TopDown baseline are \emph{directionally consistent} (positive in 14/18 configurations, mean delta $+0.0058$) but \emph{statistically fragile}: only the Qwen3-4B DL19 gain ($+0.0446$) survives Bonferroni correction. We therefore frame DualEnd as a quality-first refinement, not a free improvement.
\item The mechanism we point to is positional. Standalone worst-selection is dominated by a strong recency bias (the model labels the last passage as ``worst'' 40--63\% of the time across models), while the same model placed in a joint best-and-worst prompt produces a \emph{reversed} primacy-heavy ``dual\_worst'' pattern. Joint elicitation therefore changes the model's behavior, not just its output channel.
\end{itemize}

The paper's central scientific contribution is therefore not ``DualEnd is universally faster or better'' but \emph{worst-selection is not the dual of best-selection, and the gap between standalone worst and joint best-worst is mechanistically meaningful}. The most useful practical contribution is the empirical map: which sorting algorithm and which selection direction works for which model family, and where the joint signal is worth its compute.

\textbf{Contributions.} We summarize as follows:
\begin{enumerate}[leftmargin=*]
\item We propose three novel bidirectional strategies for LLM-based setwise ranking that extract more information per LLM call (\S\ref{sec:method}).
\item We introduce two sorting algorithms---cocktail shaker sort and double-ended selection sort---designed for dual-end LLM comparisons (\S\ref{sec:algorithms}).
\item We provide formal complexity and information-theoretic analysis comparing all methods (\S\ref{sec:algorithms}).
\item We conduct comprehensive experiments on TREC DL 2019/2020 with nine LLMs across three architecture families (Flan-T5, Qwen3, and Qwen3.5, 780M--27B), including paired significance testing and a quality--cost Pareto analysis. We complement these with a position-bias study, ranking-agreement analysis, query-difficulty stratification, and a ``when DualEnd helps'' qualitative summary on representative models (\S\ref{sec:results}, \S\ref{sec:analysis}).
\item We outline a refinement package (Selective DualEnd, bias-aware DualEnd, same-call worst-signal regularization) implemented as runnable code paths and motivated by our analyses (\S\ref{sec:analysis}).
\item We provide an open-source implementation integrated into the \texttt{llm-rankers} framework\footnote{Code and run files: \url{https://anonymous}}.
\end{enumerate}
